% ============================================================
%  FICHE REX (Retour d'Expérience) — Phase 5 : Clôture
%  Time'Eats — Cellule PMO
% ============================================================
\documentclass[11pt,a4paper,oneside]{report}
\usepackage{../style/consulting}
\hypersetup{pdftitle={Fiche REX — Time'Eats PMO}}
\pagestyle{fancy}\fancyhf{}
\renewcommand{\headrulewidth}{0pt}
\fancyhead[L]{\begin{tikzpicture}[remember picture,overlay]
  \draw[cnNavy,line width=0.5pt](0,0)--(\textwidth,0);
\end{tikzpicture}\vspace{2pt}\timeeatslogo\hspace{4pt}\small\textcolor{cnSlate}{Time'Eats \textbf{·} Fiche REX — \textit{[Nom projet]}}}
\fancyhead[R]{\small\textcolor{cnSlate}{\thepage}}
\fancyfoot[C]{\tiny\textcolor{cnSlate}{Document confidentiel — Cellule PMO — CESI MAALSI 2026}}
\begin{document}\small

\begin{center}
{\LARGE\bfseries\color{cnNavy} Fiche de Retour d'Expérience (REX)}\\[4pt]
\textcolor{cnOrange}{\rule{10cm}{1pt}}\\[4pt]
{\large\color{cnSlate} Phase 5 — Clôture \quad|\quad Projet : \textit{[Nom du projet]}}
\end{center}

\vspace{4pt}
\begin{decision}
Cette fiche est \textbf{archivée dans la base de connaissances PMO} (SharePoint). Elle alimente la revue semestrielle du référentiel. Sa rédaction est \textbf{obligatoire} pour tout projet clôturé.
\end{decision}

\vspace{4pt}
\renewcommand{\arraystretch}{1.4}
\begin{tabularx}{\textwidth}{L{3.5cm} Y L{3.5cm} Y}
\toprule
\rowcolor{cnNavy}\th{Champ} & \th{Valeur} & \th{Champ} & \th{Valeur} \\
\midrule
Projet          & \textit{[Nom]}         & Type de projet & \textit{[CRM / Infra / Web / etc.]} \\
\rowcolor{cnIce}
Chef de projet  & \textit{[Prénom NOM]}  & Date de rédaction & \textit{[JJ/MM/AAAA]} \\
Durée réelle    & \textit{[X mois]}      & Budget réalisé & \textit{[X K€ / XXX K€ prévu]} \\
\bottomrule
\end{tabularx}

\section*{1 — Ce qui a bien fonctionné}

\renewcommand{\arraystretch}{1.3}
\begin{tabularx}{\textwidth}{C{0.6cm} L{5cm} Y C{3.5cm}}
\toprule
\rowcolor{cnNavy}\th{\#} & \th{Pratique réussie} & \th{Contexte / Preuve} & \th{À répliquer ?} \\
\midrule
1 & \textit{[ex. Implication Key Users dès la phase 2]} & \textit{[explication]} & Oui — systématiser \\
\rowcolor{cnIce}
2 & \textit{[ex. Contrat forfaitaire ESN bien cadré]} & \textit{[explication]} & Oui \\
3 & \textit{[ex. Rituels de gouvernance respectés]} & \textit{[explication]} & Oui \\
\bottomrule
\end{tabularx}

\section*{2 — Ce qui aurait pu être amélioré}

\renewcommand{\arraystretch}{1.3}
\begin{tabularx}{\textwidth}{C{0.6cm} L{5cm} Y C{3.5cm}}
\toprule
\rowcolor{cnNavy}\th{\#} & \th{Problème rencontré} & \th{Cause racine} & \th{Recommandation} \\
\midrule
1 & \textit{[ex. Spécifications initiales incomplètes]} & \textit{[Pas d'atelier métier en phase 1]} & \textit{[ex. Ajouter atelier spéc. obligatoire]} \\
\rowcolor{cnIce}
2 & \textit{[ex. Retard recette de 3 semaines]} & \textit{[Environnement de recette livré en retard par ESN]} & \textit{[ex. SLA environnement à contractualiser]} \\
3 & \textit{[ex. Résistance utilisateurs]} & \textit{[Conduite du changement lancée trop tard]} & \textit{[ex. Démarrer CdC dès la phase 2]} \\
\bottomrule
\end{tabularx}

\section*{3 — Actions d'amélioration pour le PMO}

\renewcommand{\arraystretch}{1.3}
\begin{tabularx}{\textwidth}{C{0.6cm} Y L{3.5cm} C{2.5cm}}
\toprule
\rowcolor{cnNavy}\th{\#} & \th{Action à intégrer dans le référentiel PMO} & \th{Responsable} & \th{Priorité} \\
\midrule
A1 & \textit{[ex. Mettre à jour le template PMP section spécifications]} & PMO & Haute \\
\rowcolor{cnIce}
A2 & \textit{[ex. Ajouter clause SLA environnement dans contrat-type ESN]} & DSI & Normale \\
A3 & \textit{[ex. Créer checklist conduite du changement phase 2]} & PMO & Normale \\
\bottomrule
\end{tabularx}

\section*{4 — Indicateurs finaux du projet}

\renewcommand{\arraystretch}{1.3}
\begin{tabularx}{\textwidth}{L{5cm} C{3cm} C{3cm} Y}
\toprule
\rowcolor{cnNavy}\th{Indicateur} & \th{Cible} & \th{Réalisé} & \th{Commentaire} \\
\midrule
IPC final              & $\geq$ 1.00 & \textit{[X.XX]} & \\
\rowcolor{cnIce}
IPD final              & $\geq$ 1.00 & \textit{[X.XX]} & \\
Satisfaction MOA (NPS) & $\geq$ 7/10 & \textit{[X/10]} & \\
\rowcolor{cnIce}
Taux adoption J+30     & $\geq$ 70\% & \textit{[X\%]}  & \\
\bottomrule
\end{tabularx}

\section*{5 — Capitalisation}

\begin{reco}
\textbf{Archivage :} déposer cette fiche dans SharePoint sous \texttt{PMO / REX / [Année] / [Nom projet].} \\
\textbf{Diffusion :} partager en revue semestrielle PMO pour mise à jour du référentiel.
\end{reco}

\end{document}
